\section{Algorithmes de recherche}

\subsection{Fonctions de priorité}

Les stratégies de recherche sont différenciées principalement par la fonction
utilisée pour ordonner la frontière \code{OPEN}. Le \cref{tab:priorites}
présente les fonctions de priorité implantées.

\begin{table}[H]
\centering
\caption{Fonctions de priorité des stratégies de recherche.}
\label{tab:priorites}
\begin{tabular}{lll}
\toprule
\textbf{Stratégie} & \textbf{Priorité} & \textbf{Interprétation} \\
\midrule
\BFS{} & $f(n)=\mathrm{depth}(n)$ & Explore par niveau de profondeur \\
\UCS{} & $f(n)=g(n)$ & Favorise le coût cumulé minimal \\
\Greedy{} & $f(n)=h(n)$ & Favorise l'état estimé le plus proche du but \\
\AStar{} & $f(n)=g(n)+h(n)$ & Combine coût passé et estimation restante \\
\bottomrule
\end{tabular}
\end{table}

\subsection{Breadth-First Search}

\BFS{} explore les noeuds dans l'ordre croissant de profondeur. Lorsque tous
les coûts d'action sont identiques, l'algorithme trouve une solution de
profondeur minimale. Cependant, il ne garantit pas une solution de coût minimal
dans un environnement pondéré.

\subsection{Uniform-Cost Search}

\UCS{} sélectionne le noeud avec le plus faible coût de chemin $g(n)$. Avec
des coûts d'action non négatifs, il retourne une solution de coût optimal.
Son principal inconvénient est que le nombre de noeuds explorés peut être élevé
lorsque l'espace d'états est important.

\subsection{Greedy Best-First Search}

La recherche gloutonne sélectionne le noeud ayant la plus faible valeur
heuristique $h(n)$. Elle peut trouver rapidement une solution si l'heuristique
est informative, mais elle ne garantit pas une solution optimale.

\subsection{A*}

\AStar{} utilise la fonction :

\[
f(n) = g(n) + h(n)
\]

Lorsque $h(n)$ est admissible, c'est-à-dire qu'elle ne surestime jamais le coût
réel restant, A* peut retourner une solution optimale. Une heuristique
consistante renforce ce comportement dans une recherche sur graphe.

\subsection{Tree-Search et Graph-Search}

En mode \TreeSearch{}, aucun mécanisme global n'empêche la génération répétée
d'un même état. Ce mode est simple, mais il peut être très coûteux en présence
de cycles ou de chemins convergents.

En mode \GraphSearch{}, le système conserve une structure \code{best\_g} qui
associe à chaque état le meilleur coût connu. Un noeud est ignoré si un chemin
strictement meilleur vers le même état a déjà été rencontré. Ce mécanisme réduit
les duplications et améliore généralement l'efficacité de la recherche.